\documentclass[a4paper, 12pt]{article}
\usepackage[margin=1in]{geometry}
\usepackage{graphicx}
\usepackage{amsfonts, amssymb, amsmath, amsthm}
\usepackage{tikz}
\usepackage{minted}
\usepackage{hyperref}

\newcommand{\ceil}[1]{\lceil #1 \rceil}
\newcommand{\floor}[1]{\lfloor #1 \rfloor}
\newcommand{\lt}{\left}
\newcommand{\rt}{\right}
\newcommand{\C}{\mathbb{C}}
\newcommand{\R}{\mathbb{R}}
\newcommand{\Q}{\mathbb{Q}}
\newcommand{\Z}{\mathbb{Z}}
\newcommand{\N}{\mathbb{N}}
\newcommand{\zero}{\mathbf{0}}
\newcommand{\epf}{\hfill \qed}
\newcommand{\chap}[1]{
\section*{#1}
\addcontentsline{toc}{section}{#1}
}
\newcommand{\ex}[1]{
\subsection*{#1}
\addcontentsline{toc}{subsection}{#1}
}

\begin{document}

\tableofcontents
\newpage

\chap{Chapter 1 Classical Algebra}

\ex{Exercise 1.2}

Assume for contradiction that there exists $a,b \in \Z$ with $b \neq 0$ such that $\lt(\frac{a}{b}\rt)^2=2$. If $a$ is negative and $b$ is positive, then replacing $a$ with $-a$ gives us $\lt(\frac{-a}{b}\rt)^2=\lt(\frac{a}{b}\rt)^2=2$. Therefore, $-a,b$ is also a valid pair of integers satisfying the conditions. If $b$ is negative and $a$ is positive, then replacing $b$ with $-b$ gives us $\lt(\frac{a}{-b}\rt)^2=\lt(\frac{a}{b}\rt)^2=2$. Therefore, $a,-b$ is also a valid pair of integers satisfying the conditions. If both $a$ and $b$ are negative, then replacing $a$ with $-a$ and $b$ with $-b$ gives us $\lt(\frac{-a}{-b}\rt)^2=\lt(\frac{a}{b}\rt)^2=2$. Therefore, $-a,-b$ is also a valid pair of integers satisfying the conditions. Hence, we can assume without loss of generality that $a,b>0$.

Since we assume that $b$ only takes values in $\N$, by the well-ordering principle, there exists a smallest $b$ satisfying the conditions. Note that
\begin{align*}
    \lt(\frac{2b-a}{a-b}\rt)^2 &= \frac{4b^2-4ab+a^2}{a^2-2ab+b^2} \\
    &= \frac{4-4\lt(\frac{a}{b}\rt)+\lt(\frac{a}{b}\rt)^2}{\lt(\frac{a}{b}\rt)^2-2\lt(\frac{a}{b}\rt)+1} \\
    &= \frac{4-4\lt(\frac{a}{b}\rt)+2}{2-2\lt(\frac{a}{b}\rt)+1} \\
    &= \frac{6-4\lt(\frac{a}{b}\rt)}{3-2\lt(\frac{a}{b}\rt)} \\
    &= \frac{6-4\lt(\frac{a}{b}\rt)}{3-2\lt(\frac{a}{b}\rt)} \times \frac{3+2\lt(\frac{a}{b}\rt)}{3+2\lt(\frac{a}{b}\rt)} \\
    &= \frac{18 + 12\lt(\frac{a}{b}\rt) - 12\lt(\frac{a}{b}\rt) - 8\lt(\frac{a}{b}\rt)^2}{9-4\lt(\frac{a}{b}\rt)^2} \\
    &= \frac{18-8\times2}{9-4\times2} \\
    &= 2.
\end{align*}
Recall that $\lt(\frac{a}{b}\rt)^2=2$, so we have $a^2=2b^2$ and thus $2b^2-a^2=0$. Note that
\begin{align*}
    (2b-a)(2b+a) &= 4b^2-a^2 \\
    &= 2b^2 + (2b^2-a^2) \\
    &= 2b^2 \\
    &> 0
\end{align*}
since $b>0$ by assumption. Since $a,b>0$, we have $2b+a>0$. Hence, $2b-a>0$. Also, from $a^2=2b^2$ we have $a^2-b^2=b^2$. Considering $b>0$, we have
\begin{align*}
    b^2 &> 0 \\
    a^2-b^2 &> 0 \\
    (a-b)(a+b) &> 0 \\
    a-b &> 0
\end{align*}
since $a+b>0$. From $2b-a>0$ and $a-b>0$ we know that $2b-a,a-b$ is also a pair of positive integers satisfying the conditions. However, from $2b-a>0$ we have $a<2b$, and hence $a-b<2b-b=b$, contradicting the minimality of $b$. \epf

\ex{Exercise 1.5}

Let $S = \{p+q\alpha+r\alpha^2 \mid p,q,r\in\Q\}$. We first prove that $S$ is a subring of $\C$. 

Note that $1=1+0\alpha+0\alpha^2\in S$. Let $p+q\alpha+r\alpha^2,p'+q'\alpha+r'\alpha^2 \in S$, where $p,q,r,p',q',r' \in \Q$. Then we have $(p+q\alpha+r\alpha^2)+(p'+q'\alpha+r'\alpha^2) = (p+p') + (q+q')\alpha + (r+r')\alpha^2 \in S$, $-(p+q\alpha+r\alpha^2) = (-p) + (-q)\alpha + (-r)\alpha^2 \in S$ and $(p+q\alpha+r\alpha^2)(p'+q'\alpha+r'\alpha^2) = (pp'+2rq') + (pq'+p'q+2rr')\alpha + (pr'+qq'+qr'+rp')\alpha^2 \in S$. By the subring test, $S$ is a ring.

Next, we construct an inverse explicitly to show that $S$ is a subfield. Let $x = p+q\alpha+r\alpha^2 \in S \setminus \{0\}$, where $p,q,r \in \Q$. Then we have $x^2=(p^2+4qr)+(2pq+2r^2)\alpha+(2pr+q^2)\alpha^2$ and $x^3=(p^3+12pqr+2q^3+4r^3)+(3p^2q+6pr^2+6q^2r)\alpha+(3p^2r+3pq^2+6qr^2)\alpha^2$. Let $x^3+k_2x^2+k_1x+k_0=0$. Denote the coefficient of $\alpha^j$ in $x^i$ by $x_{ij}$. Then we have $(x_{30}+x_{31}\alpha+x_{32}\alpha^2)+k_2(x_{20}+x_{21}\alpha+x_{22}\alpha^2)+k_1(p+q\alpha+r\alpha^2)+k_0=0$, that is, $(x_{30}+x_{20}k_2+pk_1+k_0) + (x_{31}+x_{21}k_2+qk_1)\alpha + (x_{32}+x_{22}k_2+rk_1)\alpha^2 = 0$. Since $1,\alpha,\alpha^2$ are linearly independent, comparing coefficients gives us the linear system
\[ \begin{cases}
    k_0 + pk_1 + x_{20}k_2 = -x_{30} \\
    qk_1 + x_{21}k_2 = -x_{31} \\
    rk_1 + x_{22}k_2 = -x_{32}
\end{cases}. \]
We can use the following python code to solve this system.

\begin{minted}[frame=lines, linenos, fontsize=\footnotesize]{python}
import sympy as sp

k0, k1, k2 = sp.symbols('k0 k1 k2')
p, q, r = sp.symbols('p q r')

x30 = p**3 + 12 * p * q * r + 2 * (q**3) + 4 * (r**3)
x31 = 3 * (p**2) * q + 6 * p * (r**2) + 6 * (q**2) * r
x32 = 3 * (p**2) * r + 3 * p * (q**2) + 6 * q * (r**2)
x20 = p**2 + 4 * q * r
x21 = 2 * p * q + 2 * (r**2)
x22 = 2 * p * r + q**2

M = sp.Matrix([
    [1, p, x20, - x30],
    [0, q, x21, - x31],
    [0, r, x22, - x32]
])

solution = sp.solve_linear_system(M, k0, k1, k2)

for var in [k0, k1, k2]:
    print(f"{var} =")
    sp.pprint(sp.simplify(solution[var]))
\end{minted}

The solution is $k_0 = -p^3+6pqr-2q^3-4r^3$, $k_1=3p^2-6qr$ and $k_2=-3p$. Since $x^3+k_2x^2+k_1x+k_0 = 0$, we have $x(x^2+k_2x+k_1+k_0x^{-1}) = 0$. Recall that $\C$ is a field and hence a domain, and $x \neq 0$. Therefore, $x^2+k_2x+k_1+k_0x^{-1} = 0$, that is,
\[x^{-1} = -\frac{1}{k_0}x^2 -\frac{k_2}{k_0}x -\frac{k_1}{k_0}.\]
Since $x$ and $x^2$ are linear combinations of $1,\alpha,\alpha^2$, $x^{-1}$ is also a linear combination of $1,\alpha,\alpha^2$. \epf

\chap{Chapter 2 The Fundamental Theorem of Algebra}

\ex{Exercise 2.1}

We first prove the existence of such an expression. If $p(t)$ is constant, then set $p(t)=q(t)$ and this is the required form, since a constant polynomial $q(t)$ has no zeros in $\Q[t]$. If $p(t)$ is not constant, then $\partial p \geq 1$. By Proposition 2.7, there exist $\alpha_1,\ldots,\alpha_n \in \C$ and $0 \neq k \in \C$ such that $p(t) = k(t-\alpha_1)\ldots(t-\alpha_n)$. The coefficient of $t^n$ in the expression on the right hand side is $k$. Since $p(t) \in \Q[t]$, we have $k \in \Q$. Assume without loss of generality that $\alpha_1,\ldots,\alpha_r \in \Q$ and $\alpha_{r+1},\ldots,\alpha_n \notin \Q$ for some $0 \leq r \leq n$. Let $q(t) = k(t-\alpha_{r+1})\ldots(t-\alpha_n)$. Then $p(t) = (t-\alpha_1)\ldots(t-\alpha_r)q(t)$. Since $\alpha_{r+1},\ldots,\alpha_n$ are the only complex roots of $q(t)$, $q(t)$ has no roots in $\Q[t]$.

Suppose \[p(t) = (t-\alpha_1)\ldots(t-\alpha_r)q(t) = (t-\hat{\alpha}_1)\ldots(t-\hat{\alpha}_s)\hat{q}(t)\] such that $\alpha_1,\ldots,\alpha_r,\hat{\alpha}_1,\ldots,\hat{\alpha}_s \in \Q$ and $q(t),\hat{q}(t)$ have no zeros in $\Q[t]$. Assume without loss of generality that $\alpha_1 \leq \ldots \leq \alpha_r$. Substitute $t = \alpha_1$, we get \[0 = (\alpha_1-\alpha_1)\ldots(\alpha_1-\alpha_r)q(\alpha_1) = (\alpha_1-\hat{\alpha}_1)\ldots(\alpha_1-\hat{\alpha}_s)\hat{q}(\alpha_1).\] Since $\hat{q}(\alpha_1) \neq 0$, we have $\alpha_1-\hat{\alpha}_1 = 0$ or $\alpha_1-\hat{\alpha}_2 = 0$ or $\ldots$ or $\alpha_1-\hat{\alpha}_s = 0$. Hence, $\alpha_1 = \hat{\alpha}_1$ or $\alpha_1 = \hat{\alpha}_2$ or $\ldots$ or $\alpha_1 = \hat{\alpha}_s$. Similarly, by substituting $t = \alpha_2, \ldots, \alpha_r$, we have $\alpha_1,\alpha_2,\ldots,\alpha_r \in \{\hat{\alpha}_1, \hat{\alpha}_2, \ldots, \hat{\alpha}_s\}$. Assume without loss of generality that $\alpha_1=\hat{\alpha}_1$, $\alpha_2=\hat{\alpha}_2$, $\ldots$, $\alpha_r=\hat{\alpha}_r$. Assume for contradiction that $s>r$, so $\hat{\alpha}_{r+1},\ldots,\hat{\alpha}_s$ exist. Since $(t-\alpha_1)\ldots(t-\alpha_r) = (t-\hat{\alpha}_1)\ldots(t-\hat{\alpha}_r) \neq \zero$, we have $q(t) = (t-\hat{\alpha}_{r+1})\ldots(t-\hat{\alpha}_s)\hat{q}(t)$. Note that $\hat{\alpha}_{r+1}$ is a rational root of the polynomial on the right hand side, contradicting that $q(t)$ has no rational roots. Hence, $\hat{\alpha}_{r+1},\ldots,\hat{\alpha}_s$ do not exist and thus $q(t) = \hat{q}(t)$. This proves the uniqueness of such an expression.

Note that $p(\alpha_j) = (\alpha_j-\alpha_1)\ldots(\alpha_j-\alpha_j)\ldots(\alpha_j-\alpha_r)q(\alpha_j) = 0$, so $\alpha_j$ is a zero of $p(t)$ in $\Q$. Conversely, we want to prove that if $x \in \Q$ and $p(x)=0$ then $x = \alpha_j$ for some $j$. Assume for contradiction that there exists $x \in \Q$ such that $p(x)=0$ and $x \neq \alpha_j$ for all $j$. Then we have $x-\alpha_j \neq 0$ for all $j$. Since $q(t)$ has no zeros in $\Q[t]$, we have $q(x) \neq 0$. Since $\Q$ is a domain, we have $p(x) = (x-\alpha_1)\ldots(x-\alpha_r)q(x) \neq 0$, which contradicts that $x$ is a root of $p(t)$. Therefore, the $\alpha_j$ are precisely the zeros of $p(t)$ in $\Q$. \epf

\ex{Exercise 2.6}

Let $f(t)$ be a polynomial with $\partial f = n$. Suppose $f(t) = a_0 + a_1 t + a_2 t^2 + \ldots + a_n t^n = 0$ for all $t \in \C$. Since $f(1)=0$, we have \[ a_0 + a_1 + a_2 + \ldots + a_n = 0. \] Since $f(2)=0$, we have \[ a_0 + 2a_1 + 2^2a_2 + \ldots + 2^na_n = 0. \] Similarly, since $f(k)=0$ for any $1 \leq k \leq n+1$, we have \[ a_0 + ka_1 + k^2a_2 + \ldots + k^na_n = 0. \] Hence, we have the linear system
\[
\begin{pmatrix}
    1 & 1 & 1 & 1 & \cdots & 1 \\
    1 & 2 & 2^2 & 2^3 & \cdots & 2^n \\
    1 & 3 & 3^2 & 3^3 & \cdots & 3^n \\
    \vdots & \vdots & \vdots & \vdots & \ddots & \vdots \\
    1 & n+1 & (n+1)^2 & (n+1)^3 & \cdots & (n+1)^n
\end{pmatrix}
\begin{pmatrix} a_0 \\ a_1 \\ a_2 \\ \vdots \\ a_n \end{pmatrix} = \begin{pmatrix} 0 \\ 0 \\ 0 \\ \vdots \\ 0 \end{pmatrix}.
\]
Note that
\[
\begin{vmatrix}
    1 & 1 & 1 & 1 & \cdots & 1 \\
    1 & 2 & 2^2 & 2^3 & \cdots & 2^n \\
    1 & 3 & 3^2 & 3^3 & \cdots & 3^n \\
    \vdots & \vdots & \vdots & \vdots & \ddots & \vdots \\
    1 & n+1 & (n+1)^2 & (n+1)^3 & \cdots & (n+1)^n
\end{vmatrix}
=
\begin{vmatrix}
    1 & 1 & 1 & \cdots & 1 \\
    1 & 2 & 3 & \cdots & n+1 \\
    1^2 & 2^2 & 3^2 & \cdots & (n+1)^2 \\
    1^3 & 2^3 & 3^3 & \cdots & (n+1)^3 \\
    \vdots & \vdots & \vdots & \ddots & \vdots \\
    1^n & 2^n & 3^n & \cdots & (n+1)^n
\end{vmatrix}
,\]
which is the Vandermonde determinant with $z_1=1$, $z_2=2$, $\ldots$, $z_{n+1}=n+1$. By Exercise 2.5, this determinant is non-zero. So the linear system has no non-trivial solutions. Hence, $a_0 = a_1 = a_2 = \ldots = a_n = 0$. Therefore, all the coefficients of $f$ are zero. \epf





\end{document}
